\subsection{可度量化空间的商空间}

若 X 是可度量化空间，R 是 X 上的一个等价关系，则商空间 X/R 未必可度量化（即使 X 局部紧 * 且 X/R 正规 * 也是如此）。然而：

命题 17. 紧可度量化空间的每个豪斯多夫商空间都是紧且可度量化的。

等价地，若 \(f\) 是紧可度量化空间 \(\mathbf{X}\) 到豪斯多夫空间 \(\mathbf{X}'\) 中的连续映射，则 \(f(\mathbf{X})\) 是 \(\mathbf{X}'\) 的可度量化子空间（第 I 章，§ 9，第 4 号，定理 2 的推论 4）。

设 \(\mathbf{X}\) 是紧可度量化空间，\(\mathbf{R}\) 是 \(\mathbf{X}\) 上的一个使 \(\mathbf{X} / \mathbf{R}\) 为豪斯多夫空间的等价关系。则 \(\mathbf{X} / \mathbf{R}\) 是紧的（第 I 章，§ 9，第 4 号，定理 2），故由命题 16，只需证明 \(\mathbf{X} / \mathbf{R}\) 的拓扑具有可数基。为此，我们要用到如下事实：\(\mathbf{R}\) 是闭的（第 I 章，§ 10，第 4 号，命题 8），并且模 \(\mathbf{R}\) 的各个等价类都是紧的。设 \(\varphi\) 是 \(\mathbf{X}\) 到 \(\mathbf{X} / \mathbf{R}\) 上的典范映射，\((\mathrm{U}_n)\) 是 \(\mathbf{X}\) 的拓扑的一个可数基。设 \(z\) 是 \(\mathbf{X} / \mathbf{R}\) 的任意一点，\(\mathbf{V}\) 是 \(z\) 在 \(\mathbf{X} / \mathbf{R}\) 中的一个邻域；则 \(\overline{\varphi}^1 (\mathbf{V})\) 是 \(\mathbf{X}\) 中紧集 \(\overline{\varphi}^1 (z)\) 的邻域。若 \(x\) 是 \(\overline{\varphi}^1 (z)\) 的任意一点，则存在一个集 \(\mathrm{U}_n\) ，它包含 \(x\) 且含于 \(\overline{\varphi}^1 (\mathrm{V})\) ；于是由波莱尔–勒贝格公理，存在一个有限开覆盖 \((\mathrm{U}_{n_k})_{1\leqslant k\leqslant r}\) ，它覆盖 \(\overline{\varphi}^1 (z)\) ，且使得：若以 \(\mathbf{W}\) 表示 \(\bigcup_{k}\mathrm{U}_{n_k},\mathrm{W}\) ，则它是 \(\overline{\varphi}^1 (z)\) 的一个含于 \(\overline{\varphi}^1 (\mathrm{V})\) 的邻域。由于 \(\mathbf{R}\) 是闭的，由此推出 \(\varphi (\mathbf{W})\) 是 \(z\) 在 \(\mathbf{X} / \mathbf{R}\) 中的一个含于 \(\mathbf{V}\) 的邻域（第 I 章，§ 5，第 4 号，命题 10）。设 \(\mathfrak{B}\) 表示形如 \(\varphi (\mathbf{W})\) 的集的内部所成的集，其中 \(\mathbf{W}\) 取遍集 \(\mathfrak{F}\) ，即形如 \(U_{n}\) 的集的一切有限并所成的集；这样我们便证明了 B 是 X/R 的拓扑的一个基，又因 F 是可数的，故 B 也是可数的。

命题 18. 设 X 是完备度量空间，R 是 X 上使 X/R 为豪斯多夫空间的开等价关系，\(\varphi: X \to X/R\) 是典范映射。则对 X/R 的任一紧子集 K，存在 X 的一个紧子集 \(K'\) 使 \(\varphi(K') = K\) 。

设 \(\mathfrak{B}_1\) 是 X 中以 \(1/2\) 为半径的全体开球所成的集。当 B 取遍 \(\mathfrak{B}_1\) 时，诸集 \(\varphi(B)\) 构成 K 的一个开覆盖，因此存在 X 的有限多个点 \(x_1, \ldots, x_m\) ，使得 \(\varphi\) 把以 \(1/2\) 为半径、以 \(x_i\) （ \(1 \leqslant i \leqslant m\) ）为中心的诸开球映为构成 K 的一个开覆盖的诸集。令 \(H_1 = \{x_1, \ldots, x_m\}\) ，并设已定义了有限集 \(H_i\) （ \(1 < i \leqslant n\) ），使得：

\begin{enumerate}
\def\labelenumi{(\roman{enumi})}
\item
  \(\mathbf{H}_i\subset \mathbf{H}_{i + 1}\) ，且 \(\mathbf{H}_{i + 1}\) 的每一点都以 \(<  \mathrm{I} / 2^{i}\) 的距离靠近 \(\mathbf{H}_i\) ，这里 \(\mathbf{I}\leqslant i\leqslant n - \mathbf{I};\)
\item
  \(\varphi\) 把以 \(1/2^i\) 为半径、以 \(\mathbf{H}_i\) 的点为中心的诸开球映为构成 \(\mathbf{K}\) 的开覆盖的诸开集，这里 \(1 \leqslant i \leqslant n\) 。
\end{enumerate}

设 \(\mathfrak{B}_{n+1}\) 是以 \(1/2^{n+1}\) 为半径、中心 \(x\) 满足 \(d(x, H_n) < 1/2^n\) （ \(d\) 为 X 上的度量）的全体开球所成的集。\(H_n\) 的性质表明，诸集 \(\varphi(B)\) （ \(B \in \mathfrak{B}_{n+1}\) ）构成 K 的开覆盖；因而存在有限集 \(L_{n+1} \subset X\) ，使得 \(\varphi\) 把以 \(1/2^{n+1}\) 为半径、中心属于 \(L_{n+1}\) 的诸开球映为构成 K 的一个开覆盖的诸集。取 \(H_{n+1} = H_n \cup L_{n+1}\) ，便知可以归纳地定义一个由 X 的有限子集组成的无限序列 \((H_n)\) ，它具有上述性质 (i) 与 (ii)。令 \(H = \bigcup_{n} H_n\) ，我们来证明 H 是预紧的。对每个 \(p > 0\) 及每点 \(z_{n+p} \in H_{n+p}\) ，存在点列 \(z_{n+i} \in H_{n+i}\) （ \(0 \leqslant i \leqslant p - 1\) ）使得

\[
d \left(z _ {n + i}, z _ {n + i + 1}\right) <   \mathrm{I} / 2 ^ {n + i} \quad \text { for } \quad \mathrm{o} \leqslant i \leqslant p - \mathrm{I};
\]

由此推出 \(d(z_n, z_{n+p}) \leqslant \sum_{i=0}^{p-1} \mathrm{I}/2^{n+i} \leqslant \mathrm{I}/2^{n-1}\) ，从而 \(d(y, \mathbf{H}_n) \leqslant \mathrm{I}/2^{n-1}\) 对一切 \(y \in \mathbf{H}\) 成立，这就证明了我们的断言。由于 \(\mathbf{X}\) 完备，\(\overline{\mathbf{H}}\) 是紧的，故 \(\varphi(\overline{\mathbf{H}})\) 是紧的。其次，我们证明 \(\mathbf{K} \subset \varphi(\overline{\mathbf{H}})\) 。若 \(z \in \mathbf{K}\) ，则按定义有 \(d(\overline{\mathbf{H}}_n, \overline{\varphi}^1(z)) \leqslant \mathrm{I}/2^n\) （对一切 \(n\) ），因而 \(d(\overline{\mathbf{H}}, \overline{\varphi}^1(z)) = 0\) ；但 \(\overline{\varphi}^1(z)\) 是闭集而 \(\overline{\mathbf{H}}\) 是紧集，故由此推出 \(\overline{\mathbf{H}} \cap \overline{\varphi}^1(z) \neq \emptyset\) （第 2 号，命题 3 后的注）；断言得证。于是，若 \(\mathbf{K}' = \overline{\mathbf{H}} \cap \overline{\varphi}^1(\mathbf{K})\) ，则 \(\mathbf{K}'\) 在 \(\overline{\mathbf{H}}\) 中闭，从而是紧的，并且由已证的结果有 \(\varphi(\mathbf{K}') = \mathbf{K}\) 。证毕。

